\documentclass[conference]{IEEEtran}

\usepackage{amsmath,amssymb,amsfonts}
\usepackage{graphicx}
\usepackage{textcomp}
\usepackage{booktabs}
\usepackage{multirow}
\usepackage{hyperref}
\usepackage{algorithm}
\usepackage{algorithmic}
\usepackage{xcolor}
\usepackage{subcaption}

\graphicspath{{../figures/}}

\newcommand{\ours}{TH-GNN}
\newcommand{\eg}{\emph{e.g.}}
\newcommand{\ie}{\emph{i.e.}}
\newcommand{\etal}{\emph{et al.}}

\begin{document}

\title{ChainGuard: Temporal Heterogeneous Graph Neural Networks for Cryptocurrency Fraud Detection}

\author{
\IEEEauthorblockN{Shihua Yu}
\IEEEauthorblockA{
Department of Computer Science and Engineering\\
New York University Tandon School of Engineering\\
Brooklyn, NY 11201\\
\texttt{sy4886@nyu.edu}}
}

\maketitle

\begin{abstract}
Cryptocurrency fraud detection is increasingly critical as decentralized finance (DeFi) grows, yet existing methods fail to capture both temporal dynamics and heterogeneous transaction relationships in blockchain networks. We propose \ours{} (Temporal Heterogeneous Graph Neural Network), a framework that combines Relational Graph Convolution (R-GCN) with temporal self-attention over graph snapshots. Our key innovation is a temporal $k$-nearest-neighbor graph augmentation strategy that creates cross-timestep edges between nodes with similar transaction patterns, breaking the timestep isolation inherent in the Elliptic Bitcoin dataset. Through systematic ablation experiments on Elliptic (203,769 nodes, 234,355 edges), we demonstrate that heterogeneous edge modeling with temporal $k$-NN augmentation achieves AUC-ROC of 0.8678, a 12.29\% improvement over the GCN baseline (0.7449) under strict temporal evaluation. Notably, our results reveal that \emph{graph augmentation strategy matters more than model complexity}: the temporal $k$-NN edge augmentation alone accounts for the majority of performance gains, outperforming six baselines including Random Forest (0.8601), GraphSAGE (0.8624), and EvolveGCN (0.7994). Code is available at \url{https://github.com/youseihuayu-wonderful/ChainGuard-Crypto-Financial-Fraud-Detection-using-Temporal-Multi-Chain-Graph-Neural-Networks}.
\end{abstract}

\begin{IEEEkeywords}
graph neural network, fraud detection, cryptocurrency, temporal graph, heterogeneous graph
\end{IEEEkeywords}

\section{Introduction}

The rapid growth of decentralized finance (DeFi) has led to an explosion of cryptocurrency transactions, with the total value locked in DeFi protocols exceeding \$100 billion as of 2024. This growth has been accompanied by a corresponding increase in fraud, with cross-chain bridge exploits alone causing over \$2 billion in losses between 2021 and 2023~\cite{chainalysis2024}. Detecting illicit transactions in blockchain networks is therefore a critical challenge for maintaining trust in the cryptocurrency ecosystem.

Traditional fraud detection methods rely on hand-crafted features and tabular classifiers~\cite{weber2019anti}, ignoring the rich graph structure of blockchain transaction networks. Recent work has applied Graph Neural Networks (GNNs) to this problem~\cite{pareja2020evolvegcn, weber2019anti}, but existing approaches face two key limitations:

\begin{enumerate}
    \item \textbf{Temporal isolation}: Real-world blockchain datasets such as Elliptic~\cite{weber2019anti} organize transactions into temporal snapshots (timesteps), where edges exist only within the same timestep. Standard GNNs cannot propagate information across timesteps, limiting their ability to capture evolving fraud patterns.

    \item \textbf{Homogeneous edge treatment}: All edges are treated identically, despite representing fundamentally different types of relationships (\eg, direct payments vs.\ cross-temporal similarity patterns).
\end{enumerate}

We propose \textbf{ChainGuard}, a framework built on the Temporal Heterogeneous Graph Neural Network (\ours{}), which addresses both limitations. Our approach makes the following contributions:

\begin{itemize}
    \item \textbf{Temporal $k$-NN graph augmentation}: We introduce cross-timestep edges based on feature-space $k$-nearest-neighbor similarity between adjacent timesteps, breaking the temporal isolation of the original graph. This creates a heterogeneous graph with two edge types: original payment edges and temporal similarity edges.

    \item \textbf{Heterogeneous message passing via R-GCN}: We employ Relational Graph Convolutional Networks~\cite{schlichtkrull2018modeling} to learn separate weight matrices for each edge type, enabling the model to distinguish between direct payment relationships and temporal similarity patterns.

    \item \textbf{Comprehensive ablation study}: Through systematic ablation experiments (M1--M5), we demonstrate the individual contribution of each component and reveal a key insight: \emph{graph augmentation strategy is more important than model complexity}.

    \item \textbf{Extensive baseline comparison}: We evaluate against six baselines spanning non-graph ML (Logistic Regression, Random Forest, Gradient Boosting), standard GNNs (GAT, GraphSAGE), and temporal GNNs (EvolveGCN-H), all under strict temporal evaluation to prevent data leakage.
\end{itemize}

Our best model (M3: R-GCN with temporal $k$-NN augmentation) achieves AUC-ROC of 0.8678 on the Elliptic dataset under temporal split, outperforming all baselines including strong non-graph methods like Random Forest (0.8601) and GraphSAGE (0.8624).

\section{Related Work}

\subsection{Blockchain Fraud Detection}

Weber~\etal~\cite{weber2019anti} introduced the Elliptic dataset, the first large-scale labeled Bitcoin transaction graph, and established baselines using Logistic Regression, Random Forest, and GCN. Their work demonstrated that graph-based features can improve fraud detection over node features alone. Subsequent work~\cite{alarab2020comparative} explored additional classifiers but primarily used random train/test splits, which can leak temporal information.

\subsection{Graph Neural Networks for Financial Fraud}

GCN~\cite{kipf2017semi} has been widely applied to node classification in financial networks. GAT~\cite{velickovic2018graph} introduces attention mechanisms to weight neighbor contributions differently, while GraphSAGE~\cite{hamilton2017inductive} enables inductive learning through sampling and aggregation. However, these methods treat all edges homogeneously and cannot distinguish between different types of financial relationships.

R-GCN~\cite{schlichtkrull2018modeling} extends GCN to handle multiple edge types by learning separate weight matrices per relation, making it suitable for heterogeneous financial graphs where different edge types carry different semantic meaning.

\subsection{Temporal Graph Learning}

EvolveGCN~\cite{pareja2020evolvegcn} uses recurrent neural networks to evolve GCN weight matrices across timesteps, enabling adaptation to temporal dynamics. DySAT~\cite{sankar2020dysat} applies self-attention over graph snapshots to capture structural evolution. TGAT~\cite{xu2020inductive} combines temporal encoding with graph attention for continuous-time dynamic graphs.

These methods address temporal dynamics but typically assume cross-timestep edges already exist. In the Elliptic dataset, timesteps form completely isolated subgraphs with zero cross-timestep edges, rendering temporal GNN methods less effective without graph augmentation.

\subsection{Cross-Chain Analysis}

Recent work on cross-chain bridge security~\cite{zhang2023cross} has highlighted the need for multi-chain fraud detection. While our current evaluation uses the single-chain Elliptic dataset, our framework's heterogeneous edge modeling naturally extends to cross-chain scenarios where bridge transactions form a distinct edge type.

\section{Method: TH-GNN}

\subsection{Problem Formulation}

Let $\mathcal{G} = (\mathcal{V}, \mathcal{E}, \mathbf{X}, \mathbf{t})$ be a temporal transaction graph where $\mathcal{V}$ is the set of transaction nodes, $\mathcal{E}$ is the set of edges, $\mathbf{X} \in \mathbb{R}^{|\mathcal{V}| \times d}$ is the node feature matrix ($d = 166$ for Elliptic), and $\mathbf{t}: \mathcal{V} \rightarrow \{1, \ldots, T\}$ maps each node to its timestep. The goal is to classify each node $v \in \mathcal{V}$ as illicit ($y_v = 1$) or licit ($y_v = 0$).

\subsection{Temporal $k$-NN Graph Augmentation}
\label{sec:temporal_knn}

A key challenge in the Elliptic dataset is that each timestep forms a \emph{completely isolated subgraph}---there are zero edges between nodes in different timesteps. This prevents standard GNNs from capturing temporal patterns across time windows.

We address this by constructing temporal $k$-nearest-neighbor edges between adjacent timesteps. For each node $v$ in timestep $t$, we compute cosine similarity with all nodes in timestep $t-1$:

\begin{equation}
    \text{sim}(v, u) = \frac{\mathbf{x}_v \cdot \mathbf{x}_u}{\|\mathbf{x}_v\| \|\mathbf{x}_u\|}, \quad u \in \mathcal{V}_{t-1}
\end{equation}

and connect $v$ to its top-$k$ most similar nodes in the previous timestep ($k = 5$ in our experiments). These temporal edges are bidirectional, resulting in an augmented graph $\mathcal{G}'$ with two edge types:

\begin{itemize}
    \item \textbf{Type 0 (original)}: Direct payment/transaction edges within a timestep, representing actual fund flows.
    \item \textbf{Type 1 (temporal)}: Cross-timestep feature-similarity edges, representing temporal continuity of transaction patterns.
\end{itemize}

This augmentation adds 1,958,890 temporal edges to the original 234,355 edges, creating a connected heterogeneous graph that enables both spatial and temporal message passing.

\subsection{Heterogeneous Message Passing (R-GCN)}

Given the augmented heterogeneous graph, we employ R-GCN~\cite{schlichtkrull2018modeling} to learn separate weight matrices for each edge type. For a node $v$, the message passing at layer $l$ is:

\begin{equation}
    \mathbf{h}_v^{(l+1)} = \sigma \left( \sum_{r \in \mathcal{R}} \sum_{u \in \mathcal{N}_r(v)} \frac{1}{|\mathcal{N}_r(v)|} \mathbf{W}_r^{(l)} \mathbf{h}_u^{(l)} + \mathbf{W}_0^{(l)} \mathbf{h}_v^{(l)} \right)
\end{equation}

where $\mathcal{R} = \{0, 1\}$ denotes the two edge types, $\mathcal{N}_r(v)$ is the set of neighbors of $v$ connected by edge type $r$, $\mathbf{W}_r^{(l)}$ is the learnable weight matrix for relation $r$ at layer $l$, $\mathbf{W}_0^{(l)}$ is a self-loop weight, and $\sigma$ is the ReLU activation function.

We use a 2-layer R-GCN with hidden dimension 80 (chosen for parameter parity with the GCN baseline), followed by a linear classifier:

\begin{equation}
    \hat{y}_v = \sigma(\mathbf{w}^T \mathbf{h}_v^{(2)} + b)
\end{equation}

where $\sigma$ is the sigmoid function.

\subsection{Temporal Self-Attention (M4 Extension)}

For the full TH-GNN model (M4 in our ablation), we add a temporal self-attention layer on top of the R-GCN embeddings. We pool node embeddings by timestep:

\begin{equation}
    \mathbf{z}_t = \frac{1}{|\mathcal{V}_t|} \sum_{v \in \mathcal{V}_t} \mathbf{h}_v^{(2)}
\end{equation}

and apply multi-head causal self-attention with sinusoidal positional encoding across the timestep sequence $(\mathbf{z}_1, \ldots, \mathbf{z}_T)$. A learnable gating mechanism combines the temporal context with individual node embeddings:

\begin{equation}
    \mathbf{h}_v' = g \odot \mathbf{h}_v^{(2)} + (1 - g) \odot \mathbf{c}_{t(v)}
\end{equation}

where $g = \sigma(\mathbf{W}_g [\mathbf{h}_v^{(2)} \| \mathbf{c}_{t(v)}])$ and $\mathbf{c}_{t(v)}$ is the attention-weighted temporal context for timestep $t(v)$.

\subsection{Semi-Supervised Label Propagation (M5 Extension)}

To leverage the large number of unlabeled nodes (157,205 out of 203,769 in Elliptic), we add a label propagation module that iteratively diffuses known labels through the augmented graph:

\begin{equation}
    \mathbf{Y}^{(t+1)} = (1 - \alpha) \mathbf{Y}^{(0)} + \alpha \tilde{\mathbf{A}} \mathbf{Y}^{(t)}
\end{equation}

where $\tilde{\mathbf{A}}$ is the symmetrically normalized adjacency matrix and $\alpha = 0.5$ controls the diffusion rate. The resulting soft labels provide a consistency regularization term:

\begin{equation}
    \mathcal{L} = \mathcal{L}_{\text{BCE}} + \lambda \cdot \text{MSE}(\hat{\mathbf{y}}_U, \mathbf{y}_U^{\text{LP}})
\end{equation}

where $\hat{\mathbf{y}}_U$ are GNN predictions on unlabeled nodes and $\mathbf{y}_U^{\text{LP}}$ are label propagation soft labels ($\lambda = 0.1$).

\section{Experiments}

\subsection{Dataset}

We evaluate on the \textbf{Elliptic Bitcoin Dataset}~\cite{weber2019anti}, the largest publicly available labeled cryptocurrency transaction graph. Table~\ref{tab:dataset} summarizes its statistics.

\begin{table}[h]
\centering
\caption{Elliptic Dataset Statistics}
\label{tab:dataset}
\begin{tabular}{lr}
\toprule
\textbf{Property} & \textbf{Value} \\
\midrule
Nodes (transactions) & 203,769 \\
Edges (payment flows) & 234,355 \\
Temporal edges ($k$-NN, $k=5$) & 1,958,890 \\
Node features & 166 \\
Timesteps & 49 \\
Illicit (labeled) & 4,545 (9.8\%) \\
Licit (labeled) & 42,019 (90.2\%) \\
Unknown (unlabeled) & 157,205 \\
\bottomrule
\end{tabular}
\end{table}

\subsection{Experimental Setup}

\textbf{Temporal split.} We use a strict temporal train/validation/test split to prevent data leakage: timesteps 1--34 for training (29,894 labeled nodes), 35--41 for validation (7,829), and 42--49 for testing (8,841). This is more realistic than random splitting, which can leak future information.

\textbf{Evaluation metrics.} We report AUC-ROC (primary metric), F1 score, Precision, and Recall for the illicit class. All experiments use seed 42 for reproducibility.

\textbf{Training details.} All GNN models are trained with Adam optimizer (lr=0.01, weight decay=$5 \times 10^{-4}$), binary cross-entropy loss with class weight 7.63 (licit/illicit ratio), dropout 0.5, and early stopping with patience 20 based on validation AUC.

\subsection{Ablation Study}

We conduct a systematic ablation study by incrementally adding components to the GCN baseline (Table~\ref{tab:ablation}).

\begin{table}[h]
\centering
\caption{Ablation Study Results (Temporal Split)}
\label{tab:ablation}
\begin{tabular}{llcccc}
\toprule
\textbf{ID} & \textbf{Model} & \textbf{AUC} & \textbf{F1} & \textbf{Prec.} & \textbf{Recall} \\
\midrule
M1 & GCN & 0.7449 & 0.281 & 0.278 & 0.284 \\
M2 & +Temporal Attn & 0.7937 & 0.366 & 0.461 & 0.304 \\
\textbf{M3} & \textbf{+Hetero (R-GCN)} & \textbf{0.8678} & \textbf{0.511} & \textbf{0.717} & 0.397 \\
M4 & Full TH-GNN & 0.8535 & 0.493 & 0.613 & 0.412 \\
M5 & +Label Prop. & 0.8435 & 0.474 & 0.659 & 0.370 \\
\bottomrule
\end{tabular}
\end{table}

\textbf{Key findings from ablation:}

\begin{enumerate}
    \item \textbf{M3 provides the largest single-step improvement} (+12.29\% AUC over M1). The heterogeneous edge modeling with temporal $k$-NN augmentation is the dominant contributor to performance gains.

    \item \textbf{M4 $<$ M3}: Adding temporal self-attention on top of R-GCN with temporal edges provides diminishing returns (AUC 0.8535 vs.\ 0.8678). This is because temporal $k$-NN edges already encode temporal patterns, making explicit temporal attention partially redundant.

    \item \textbf{Label propagation (M5) does not further improve}: The consistency regularization adds noise rather than signal, suggesting that the GNN has already learned effective representations from the augmented graph.

    \item \textbf{Graph augmentation $>$ model complexity}: The simplest model with the right graph structure (M3) outperforms more complex architectures (M4, M5) on the same task.
\end{enumerate}

\subsection{Baseline Comparison}

We compare against six baselines spanning three categories (Table~\ref{tab:baselines}).

\begin{table}[h]
\centering
\caption{Baseline Comparison Results (Temporal Split)}
\label{tab:baselines}
\begin{tabular}{llcccc}
\toprule
\textbf{Method} & \textbf{Type} & \textbf{AUC} & \textbf{F1} & \textbf{Prec.} & \textbf{Recall} \\
\midrule
LR & Non-graph & 0.855 & 0.216 & 0.126 & 0.765 \\
Random Forest & Non-graph & 0.860 & 0.620 & 0.969 & 0.456 \\
Grad.\ Boosting & Non-graph & 0.843 & 0.546 & 0.615 & 0.490 \\
\midrule
GCN (M1) & GNN & 0.745 & 0.281 & 0.278 & 0.284 \\
GAT & GNN & 0.805 & 0.288 & 0.208 & 0.463 \\
GraphSAGE & GNN & 0.862 & 0.540 & 0.751 & 0.422 \\
\midrule
EvolveGCN-H & Temp.\ GNN & 0.799 & 0.193 & 0.119 & 0.522 \\
\midrule
\textbf{\ours{} (Ours)} & \textbf{TH-GNN} & \textbf{0.868} & \textbf{0.511} & 0.717 & 0.397 \\
\bottomrule
\end{tabular}
\end{table}

\textbf{Key observations:}

\begin{enumerate}
    \item \textbf{Non-graph methods outperform standard GNNs on the original graph.} LR (0.855), RF (0.860), and Gradient Boosting (0.843) all exceed GCN (0.745) and GAT (0.805). This confirms that the original Elliptic graph---with isolated timestep subgraphs---provides limited structural information.

    \item \textbf{GraphSAGE (0.862) is the strongest standard GNN baseline}, approaching non-graph methods. Its sampling-based aggregation may be more robust to the sparse, isolated structure.

    \item \textbf{EvolveGCN-H (0.799) underperforms}, despite being designed for temporal graphs. Its GRU-based weight evolution struggles with completely isolated timesteps where each snapshot is an independent subgraph.

    \item \textbf{\ours{} (0.868) achieves the highest AUC-ROC}, outperforming all baselines. The temporal $k$-NN augmentation transforms the graph into a connected structure where GNN message passing is effective, giving our method a decisive advantage over both feature-only and graph-only approaches.
\end{enumerate}

\section{Case Study and Analysis}

\subsection{Detection Improvement Analysis}

To understand \emph{where} \ours{} improves over the GCN baseline, we compare their predictions on the 408 illicit nodes in the test set (timesteps 42--49).

Figure~\ref{fig:case_study} shows the detection breakdown. \ours{} correctly identifies a substantial number of illicit nodes that GCN misses (M3-only detections), while the predicted probability distributions show that \ours{} produces more confident predictions for illicit nodes, with a clearer separation from the decision threshold at 0.5.

\subsection{Why Graph Augmentation Matters Most}

Our ablation reveals a counterintuitive finding: the simplest heterogeneous model (M3) outperforms more complex variants (M4, M5). We attribute this to the nature of the temporal $k$-NN augmentation:

\begin{enumerate}
    \item \textbf{Breaking isolation is the bottleneck.} The original Elliptic graph has 49 completely disconnected timestep subgraphs. Standard GNNs can only aggregate within-timestep neighborhoods of 2--5 hops. By adding temporal edges, we connect the graph and enable cross-timestep information flow.

    \item \textbf{Feature similarity encodes temporal patterns.} The cosine-similarity-based $k$-NN edges implicitly capture temporal continuity: nodes with similar features across adjacent timesteps are likely part of the same evolving fraud pattern. This makes explicit temporal attention (M4) partially redundant.

    \item \textbf{Heterogeneous weights are important.} Original payment edges and temporal similarity edges carry fundamentally different information. R-GCN's per-type weight matrices (M3) outperform uniform treatment (M1 on augmented graph) by a significant margin, confirming that edge type distinction is valuable.
\end{enumerate}

\subsection{Limitations}

Our evaluation has several limitations that we plan to address in future work:

\begin{itemize}
    \item \textbf{Single dataset}: We evaluate only on Elliptic. Validation on additional blockchain datasets would strengthen our claims.
    \item \textbf{Single seed}: Results are reported for seed 42. Multi-seed runs with standard deviation would provide statistical robustness.
    \item \textbf{Feature-based $k$-NN}: Temporal edges are constructed from raw features, which may not generalize to datasets with different feature distributions.
\end{itemize}

\section{Conclusion}

We presented \ours{}, a Temporal Heterogeneous Graph Neural Network for cryptocurrency fraud detection. Through systematic ablation experiments and comprehensive baseline comparison on the Elliptic Bitcoin dataset, we demonstrated that:

\begin{enumerate}
    \item \textbf{Temporal $k$-NN graph augmentation} is the most impactful component, breaking the timestep isolation that limits standard GNN approaches on blockchain transaction graphs.
    \item \textbf{Heterogeneous edge modeling} via R-GCN provides a 12.29\% AUC improvement over the GCN baseline, achieving the best performance among all tested methods (AUC = 0.8678).
    \item \textbf{Graph augmentation strategy matters more than model complexity}---the simplest model with the right graph structure outperforms more architecturally complex variants.
\end{enumerate}

Our findings have practical implications for blockchain security: effective fraud detection requires not just sophisticated models, but careful construction of the transaction graph to capture cross-temporal relationships.

\subsection{Future Work}

Several directions for future research include:

\begin{itemize}
    \item \textbf{Cross-chain evaluation}: Extending \ours{} to multi-chain graphs with bridge transaction edges as a third heterogeneous edge type, using datasets like XChainDataGen.
    \item \textbf{Online learning}: Adapting the model for streaming blockchain data where new timesteps arrive continuously.
    \item \textbf{Learned graph augmentation}: Replacing the fixed $k$-NN strategy with learnable edge construction (\eg, via graph structure learning).
    \item \textbf{Larger-scale validation}: Evaluating on additional blockchain datasets (Ethereum, multi-chain) to test generalizability.
\end{itemize}


\bibliographystyle{IEEEtran}
\bibliography{references}

\end{document}
